% Part: computability
% Chapter: recursive-functions
% Section: trees

\documentclass[../../../include/open-logic-section]{subfiles}

\begin{document}

\olfileid{cmp}{rec}{tre}
\olsection{الأشجار}

كما رمزنا إلى المتتاليات بأعداد طبيعية، ومثلنا خصائصها وعملياتها
بعلائق ودوال عودية بدائية تجري على تلك الأعداد، يجوز أن نصنع
بالأشجار مثل ذلك. ونستعين لهذا الغرض بالمتتاليات ورموزها. والشجرة إما
عقدة منفردة، وإما عقدة تتصل بجملة من الأشجار
الجزئية؛ ويجوز في كلتا الحالتين أن تحمل العقدة لصيقة. وهذه العقدة هي \emph{جذر} الشجرة،
وما يتصل بها من الأشجار الجزئية هو \emph{أشجارها الجزئية المباشرة}.

فترميز الأشجار عودي، ورمز الشجرة متتالية $\tuple{{\OLId{latin-ordinary}{k}},{\OLId{latin-ordinary}{d}}_{\OLMathNumeral{1}},\dots,{\OLId{latin-ordinary}{d}}_{\OLId{latin-ordinary}{k}}}$، حيث ${\OLId{latin-ordinary}{k}}$ عدد
الأشجار الجزئية المباشرة، وحيث ${\OLId{latin-ordinary}{d}}_{\OLMathNumeral{1}}$، …،~${\OLId{latin-ordinary}{d}}_{\OLId{latin-ordinary}{k}}$ رموز هذه الأشجار
الجزئية. فإن كانت للعقد لصائق، ألحقناها برموز الأشجار الجزئية
المباشرة. وعلى هذا فالشجرة المؤلفة من عقدة واحدة ذات لصيقة~${\OLId{latin-ordinary}{l}}$ يرمز إليها
بـ$\tuple{{\OLMathNumeral{0}},{\OLId{latin-ordinary}{l}}}$، ونرمز إلى شجرة تتألف من جذر (ذو لصيقة~${\OLId{latin-ordinary}{l}}_{\OLMathNumeral{1}}$) متصل
بعقدتين منفردتين (ذواتي اللصيقتين ${\OLId{latin-ordinary}{l}}_{\OLMathNumeral{2}}$ و${\OLId{latin-ordinary}{l}}_{\OLMathNumeral{3}}$) بـ
$\tuple{{\OLMathNumeral{2}},\tuple{{\OLMathNumeral{0}},{\OLId{latin-ordinary}{l}}_{\OLMathNumeral{2}}},\tuple{{\OLMathNumeral{0}},{\OLId{latin-ordinary}{l}}_{\OLMathNumeral{3}}},{\OLId{latin-ordinary}{l}}_{\OLMathNumeral{1}}}$.

\begin{prop}
  \ollabel{prop:subtreeseq}
  الدالة $\fn{SubtreeSeq}({\OLId{latin-ordinary}{t}})$، التي تعيد رمز متتالية عناصرها رموز جميع
  الأشجار الجزئية للشجرة ذات الرمز~${\OLId{latin-ordinary}{t}}$، عودية بدائية.
\end{prop}

\begin{proof}
  نبدأ بأن $\fn{ISubtrees}({\OLId{latin-ordinary}{t}})=\fn{subseq}({\OLId{latin-ordinary}{t}},{\OLMathNumeral{1}},({\OLId{latin-ordinary}{t}})_{\OLMathNumeral{0}})$ عودية
  بدائية، وأنها تعيد رموز الأشجار الجزئية المباشرة لشجرة~${\OLId{latin-ordinary}{t}}$. ثم
  نعرّف دالة مساعدة $\fn{hSubtreeSeq}({\OLId{latin-ordinary}{t}},{\OLId{latin-ordinary}{n}})$ تحسب متتالية جميع
  الأشجار الجزئية التي لا يزيد بعدها عن الجذر على ${\OLId{latin-ordinary}{n}}$ عقد. أما الأشجار الجزئية
  لـ~${\OLId{latin-ordinary}{t}}$ التي بعدها ${\OLMathNumeral{0}}$ عقد عن الجذر، أعني التي تبدأ من جذر~${\OLId{latin-ordinary}{t}}$، فمتتاليتها
  مؤلفة من~${\OLId{latin-ordinary}{t}}$ وحدها. ولطلب جميع الأشجار
  الجزئية حتى المستوى~${\OLId{latin-ordinary}{n}}+{\OLMathNumeral{1}}$ من الشجرة ${\OLId{latin-ordinary}{t}}$، نصل متتالية ما بلغ المستوى
  ${\OLId{latin-ordinary}{n}}$ بمتتالية جميع الأشجار الجزئية المباشرة لما وقع في المستويات
  حتى ${\OLId{latin-ordinary}{n}}$. ولجمع هذه في قائمة واحدة نستعمل الأمر الآتي: متى كانت ${\OLId{latin-ordinary}{f}}({\OLId{latin-ordinary}{x}})$ دالة عودية
  بدائية تعيد رموز متتاليات، فإن الدالة ${\OLId{latin-ordinary}{g}}_{\OLId{latin-ordinary}{f}}({\OLId{latin-ordinary}{s}},{\OLId{latin-ordinary}{k}})$ التي تصل قيم ${\OLId{latin-ordinary}{f}}$
  لأول ${\OLId{latin-ordinary}{k}}$ من عناصر ${\OLId{latin-ordinary}{s}}$ عودية بدائية أيضًا:
    \begin{align*}
      {\OLId{latin-ordinary}{g}}_{\OLId{latin-ordinary}{f}}({\OLId{latin-ordinary}{s}}, {\OLMathNumeral{0}}) & = \emptyseq\\
      {\OLId{latin-ordinary}{g}}_{\OLId{latin-ordinary}{f}}({\OLId{latin-ordinary}{s}}, {\OLId{latin-ordinary}{k}}+{\OLMathNumeral{1}}) & = {\OLId{latin-ordinary}{g}}_{\OLId{latin-ordinary}{f}}({\OLId{latin-ordinary}{s}}, {\OLId{latin-ordinary}{k}}) \concat {\OLId{latin-ordinary}{f}}(({\OLId{latin-ordinary}{s}})_{\OLId{latin-ordinary}{k}})
    \end{align*}
    فمثلًا، إذا كانت ${\OLId{latin-ordinary}{s}}$ متتالية أشجار، فإن
    ${\OLId{latin-ordinary}{h}}({\OLId{latin-ordinary}{s}})={\OLId{latin-ordinary}{g}}_{\fn{ISubtrees}}({\OLId{latin-ordinary}{s}},\len{{\OLId{latin-ordinary}{s}}})$ تعطي متتالية الأشجار الجزئية
    المباشرة لعناصر~${\OLId{latin-ordinary}{s}}$. ومنها نحصل على تعريف
    $\fn{hSubtreeSeq}$ بـ
    \begin{align*}
      \fn{hSubtreeSeq}({\OLId{latin-ordinary}{t}}, {\OLMathNumeral{0}}) & = \tuple{{\OLId{latin-ordinary}{t}}} \\
      \fn{hSubtreeSeq}({\OLId{latin-ordinary}{t}}, {\OLId{latin-ordinary}{n}}+{\OLMathNumeral{1}}) & = \fn{hSubtreeSeq}({\OLId{latin-ordinary}{t}}, {\OLId{latin-ordinary}{n}}) \concat
      {\OLId{latin-ordinary}{h}}(\fn{hSubtreeSeq}({\OLId{latin-ordinary}{t}}, {\OLId{latin-ordinary}{n}})).
    \end{align*}
    ولا يتجاوز أقصى مستوى في الشجرة ذات الرمز~${\OLId{latin-ordinary}{t}}$، وهو أبعد مسافة بين
    الجذر وورقة، عدد الترميز~${\OLId{latin-ordinary}{t}}$. فتكون قيمة
    $\fn{hSubtreeSeq}({\OLId{latin-ordinary}{t}},{\OLId{latin-ordinary}{t}})$ متتالية رموز جميع الأشجار الجزئية للشجرة
    ذات الرمز~${\OLId{latin-ordinary}{t}}$.
\end{proof}

\begin{prob}
  قد يكرر تعريف $\fn{hSubtreeSeq}$ الوارد في برهان
  \olref[cmp][rec][tre]{prop:subtreeseq} بعض الرموز. فضع بدلًا منه تعريفًا
  لا يورد رمز أي شجرة جزئية في القائمة الناتجة إلا مرة واحدة.
\end{prob}

\end{document}
